\documentclass[12pt]{article}
\usepackage[utf8]{inputenc}
\usepackage[english]{babel}
\usepackage{geometry}
\geometry{a4paper, margin=2.5cm}
\usepackage{listings}
\usepackage{xcolor}
\usepackage{hyperref}

\definecolor{codegreen}{rgb}{0,0.6,0}
\definecolor{codegray}{rgb}{0.5,0.5,0.5}
\definecolor{codepurple}{rgb}{0.58,0,0.82}
\definecolor{backcolour}{rgb}{0.95,0.95,0.92}

\lstdefinestyle{mystyle}{
    backgroundcolor=\color{backcolour},   
    commentstyle=\color{codegreen},
    keywordstyle=\color{magenta},
    numberstyle=\tiny\color{codegray},
    stringstyle=\color{codepurple},
    basicstyle=\ttfamily\footnotesize,
    breakatwhitespace=false,         
    breaklines=true,                 
    captionpos=b,                    
    keepspaces=true,                 
    numbers=left,                    
    numbersep=5pt,                  
    showspaces=false,                
    showstringspaces=false,
    showtabs=false,                  
    tabsize=2
}

\lstset{style=mystyle}

\title{Implementation}
\author{NeoLAAC Development Team}
\date{\today}

\begin{document}

\maketitle

\section{General Introduction and Architectural Decision Process}
The implementation journey of the NeoLAAC (Freshman Support Line) digital system started with the analysis of support requirements for freshmen at the University of Beira Interior. The team's mental process focused on overcoming the inherent issues of traditional monolithic systems, such as excessive coupling and cascade failure propagation. Thus, it was decided to adopt a decentralized microservices architecture, distributing responsibilities across isolated Docker containers. Databases were logically segregated, ensuring that traffic spikes in student forums or support chats would not affect critical emergency services. Internal communication occurs via the HTTP protocol with JSON data payloads, centralized through an orchestrator gateway.

\section{Central Orchestrator and API Gateway}
This section presents the ingress data layer and centralized routing of the application. The central orchestrator assumes the responsibility of managing and forwarding incoming requests, guaranteeing the isolation of the internal network and the security of background services, as detailed below.

\subsection{orq Service}
The orq service acts as the unified API gateway of the application, intermediating communication between the frontend graphical interface and internal APIs. The mental process underlying this decision was to prevent direct public exposure of the multiple microservices within the ecosystem, reducing the attack surface to a single control point. The orq service validates routes, proxies asynchronous requests, and executes cross-logical checks involving multiple services, such as validating whether a user has permissions in a specific student organization before creating an event. The chat messaging logic includes critical term triage to detect risks swiftly.

The following code excerpt details the chat routing and triage of the orq service.

\begin{lstlisting}[language=Python]
@app.post("/chat/messages")
async def send_chat(data: dict):
    incident_id = data.get("incident_id")
    sender_email = data.get("sender_email")
    message = data.get("message", "")
    is_responder = data.get("is_responder", False)

    async with httpx.AsyncClient() as client:
        if is_responder:
            ESCALATED_INCIDENTS[incident_id] = True
            response = await client.post(f"{CHAT_SERVICE_URL}/messages", json=data)
            return response.json()

        response = await client.post(f"{CHAT_SERVICE_URL}/messages", json=data)
        
        if ESCALATED_INCIDENTS.get(incident_id, False):
            try:
                await client.post(f"{NOTIFICATION_SERVICE_URL}/notify", json={
                    "target_role": "Response-team",
                    "message": f"NEW MESSAGE: {sender_email} sent a message.",
                    "data": {"type": "chat_message", "incident_id": incident_id, "user": sender_email}
                })
            except:
                pass
            return response.json()

        CRITICAL_KEYWORDS = ["emergência", "urgência", "pânico", "médico", "acidente", "polícia", "agressão", "roubo", "fogo", "incêndio", "ferido", "hospital", "morrer", "humano", "pessoa", "assistente", "operador"]
        msg_lower = message.lower()
        if any(keyword in msg_lower for keyword in CRITICAL_KEYWORDS):
            ESCALATED_INCIDENTS[incident_id] = True
            bot_msg = "Understood. I identified an urgent request. I am transferring this chat immediately to our support team."
            await client.post(f"{CHAT_SERVICE_URL}/messages", json={
                "incident_id": incident_id,
                "sender_email": "mentor@ubi.pt",
                "message": bot_msg,
                "is_responder": True
            })
            try:
                await client.post(f"{NOTIFICATION_SERVICE_URL}/notify", json={
                    "target_role": "Response-team",
                    "message": f"CRITICAL ESCALATION: Incident #{incident_id} escalated to humans!",
                    "data": {"type": "chat_message", "incident_id": incident_id, "user": sender_email}
                })
            except:
                pass
            return {"status": "escalated_to_human"}
\end{lstlisting}

\section{Core User Services and Access Control}
This section describes the set of microservices dedicated to validating credentials, managing public profiles, and monitoring user presence and activity in NeoLAAC. These components ensure that access to the platform respect data privacy regulations, tracking the user journey in a persistent and secure manner.

\subsection{auth-service Service}
The auth-service service is responsible for credential storage, user registration, and login validation, generating JSON Web Tokens for access control. Early in the implementation, a captcha-based verification challenge was introduced to prevent automated login attempts. The subsequent review process led to its removal. Captcha token validation relies on external verification queries that require a static/fixed IP configuration. Since the local development environment used dynamic IP configurations provided by domestic ISPs, captcha requests failed repeatedly, blocking access and forcing the removal of this external service.

The code excerpt below reflects the captcha logic that was removed from the authentication service.

\begin{lstlisting}[language=Python]
# Captcha validation method that was removed
async def verify_captcha_token(captcha_response: str, remote_ip: str) -> bool:
    # Dynamic IP usage caused validation queries to fail
    payload = {
        "secret": CAPTCHA_PRIVATE_KEY,
        "response": captcha_response,
        "remoteip": remote_ip
    }
    async with httpx.AsyncClient() as client:
        res = await client.post("https://www.google.com/recaptcha/api/siteverify", data=payload)
        return res.json().get("success", False)

@app.post("/login")
async def login(user: UserAuth, db: Session = Depends(get_db)):
    # The following snippet was removed from the login flow:
    # if not await verify_captcha_token(user.captcha_token, user.remote_ip):
    #     raise HTTPException(status_code=400, detail="Invalid captcha")
    
    row = db.execute(text("SELECT id, password_hash, role FROM users WHERE email = :email"), {"email": user.email}).fetchone()
    if not row or not pwd_context.verify(user.password, row[1]):
        raise HTTPException(status_code=401, detail="Incorrect credentials")
    return {"access_token": create_access_token({"sub": user.email})}
\end{lstlisting}

\subsection{profile-service Service}
The profile-service service stores extended user profile information, including display names, phone numbers, and avatars. The team's mental process in separating profiles from authentication was to ensure data privacy compliance. In this manner, the critical credential table remains in an isolated database service, while public profile and social data can be exposed and queried flexibly by the social routes.

\subsection{logU-service Service}
The logU-service service functions as a subsystem integrated into the authentication module to audit user actions. The mental process behind it was to detect security anomalies, such as multiple failed login attempts or unauthorized privilege escalations. This service aggregates user events and pushes them directly to the centralized logging service, ensuring audit trails without compromising active transactional database performance.

\subsection{Usertrack-service Service}
The Usertrack-service service monitors user presence in real-time on the platform, recording activity timestamps and updating the online status of volunteers and responders. The mental process focused on this tracker aimed to optimize support ticket distribution. By tracking active staff presence, the orchestrator routes escalated chats to online responders, preventing queue stalls.

\section{Academic and Information Services}
This section covers the microservices that provide institutional and academic knowledge to new students. Through these modules, the application exposes official curriculum data, answers general questions automatically, and guarantees that freshmen have quick and reliable references during their transition.

\subsection{academic-service Service}
The academic-service service stores and exposes structured course and subject data. The mental process behind isolating this service was to establish a single source of academic truth. Other services, such as calendars or subject-specific forums, query this service to validate course registrations and maintain consistent department data.

\subsection{calendar-service Service}
The calendar-service service, built in Node.js with a dedicated MariaDB database, manages classes, schedules, and report deviations.

\subsection{ClassDates Integration}
To populate the database with exact class dates for a given academic semester, the calendar service integrates the ClassDates library. Developed by Professor Dr. Tiago M. C. Simões, this Python and JavaScript library computes recurrent class dates for a given academic year, weekdays, breaks, and national holidays. The team's mental process was to adapt this library to automate schedule generation in the calendar database, minimizing manual scheduling errors.

The following excerpt demonstrates the use of the ClassDates library in the calendar service.

\begin{lstlisting}[language=JavaScript]
// Integration of the ClassDates library in calendar-service
const classDates = require('./ClassDates/index.js');

app.get('/classes/generate', async (req, res) => {
    const { start_date, end_date, weekday, breaks_json } = req.query;
    const config = {
        format: "YYYY-MM-DD",
        start: start_date,
        end: end_date,
        weekday: parseInt(weekday),
        breaks: JSON.parse(breaks_json || "[]")
    };
    
    const generatedDates = classDates.generate(config);
    if (!generatedDates) {
        return res.status(400).json({ error: "Invalid configuration" });
    }
    res.json(generatedDates.map(d => d.format("YYYY-MM-DD")));
});
\end{lstlisting}

\subsection{FAQ-service Service}
The FAQ-service service exposes common questions and answers about UBI. The mental process in its design was simplicity and read performance optimization. Since FAQ data is mostly static and rarely changed, this service was optimized for minimum response latency, serving as a stable knowledge base for freshmen.

\subsection{chatbot-service Service}
The chatbot-service service processes user inputs in natural language, acting as the primary digital mentor. The design process focused on resource efficiency. Instead of deploying complex, resource-heavy language models, the chatbot normalizes user messages by stripping accents and punctuation, scores word overlap against FAQ questions, and returns the response or a fallback flag triggering human escalation.

The text normalization function of the chatbot is presented below.

\begin{lstlisting}[language=Python]
def normalize_text(text: str) -> str:
    text = text.lower()
    replacements = {
        'á': 'a', 'à': 'a', 'â': 'a', 'ã': 'a',
        'é': 'e', 'è': 'e', 'ê': 'e',
        'í': 'i', 'ì': 'i', 'î': 'i',
        'ó': 'o', 'ò': 'o', 'ô': 'o', 'õ': 'o',
        'ú': 'u', 'ù': 'u', 'û': 'u',
        'ç': 'c'
    }
    for orig, rep in replacements.items():
        text = text.replace(orig, rep)
    text = re.sub(r'[^\w\s]', '', text)
    return text.strip()
\end{lstlisting}

\subsection{page-service Service}
The page-service service organizes and serves static survival guides, code of conduct, and contacts. The design process focused on direct accessibility, ensuring that critical documents can be read instantly without database overhead or login requirements.

\section{Social Interaction and Gamification Services}
This section analyzes the social interaction and gamification dynamics integrated into the portal. The associated microservices support community building, enabling users to share news, register autonomous events, and promote the exploration of the university's physical space through dynamic photo challenges.

\subsection{challenge-service Service}
The challenge-service service handles thematic photo challenges. The design process focused on fostering campus exploration through gamification. The service processes photo uploads, generates unique UUID filenames, and limits entries to one submission per student. Administrators can evaluate submissions to declare a winner and close challenges.

The following excerpt illustrates photo uploads in the challenge-service.

\begin{lstlisting}[language=Python]
@app.post("/challenges/{challenge_id}/submit")
async def submit_photo(
    challenge_id: int,
    user_id: int = Form(...),
    user_email: str = Form(...),
    caption: str = Form(""),
    image: UploadFile = File(...),
    db: Session = Depends(get_db)
):
    chal = db.execute(text("SELECT status FROM challenges WHERE id = :id"), {"id": challenge_id}).fetchone()
    if not chal or chal[0] == "ended":
        raise HTTPException(status_code=400, detail="Inactive challenge")

    existing = db.execute(
        text("SELECT id FROM challenge_submissions WHERE challenge_id = :cid AND user_id = :uid"),
        {"cid": challenge_id, "uid": user_id}
    ).fetchone()
    if existing:
        raise HTTPException(status_code=400, detail="Submission already done")

    ext = image.filename.split(".")[-1]
    filename = f"{uuid.uuid4()}.{ext}"
    filepath = os.path.join(UPLOAD_DIR, filename)
    with open(filepath, "wb") as buffer:
        shutil.copyfileobj(image.file, buffer)
        
    photo_url = f"/api/challenge-static/uploads/{filename}"
    db.execute(
        text("INSERT INTO challenge_submissions (challenge_id, user_id, user_email, photo_url, caption) VALUES (:cid, :uid, :uemail, :purl, :caption)"),
        {"cid": challenge_id, "uid": user_id, "uemail": user_email, "purl": photo_url, "caption": caption}
    )
    db.commit()
    return {"status": "success", "photo_url": photo_url}
\end{lstlisting}

\subsection{post-service Service}
The post-service service implements the student social forum, allowing posts with text and media attachments, likes, and comments. The design process focused on data integrity, processing likes and comments in separate transactional tables with cascade deletions.

\subsection{news-service Service}
The news-service service handles formal campus announcements. The design process was to keep editorial news separate from informal student posts to ensure a reliable information channel.

\subsection{events-service Service}
The events-service service manages events promoted by student organizations. The design process focused on delegating autonomy, validating user permissions via the profile service before allowing event management.

\subsection{feed-service Service}
The feed-service service aggregates social posts and news into a chronological feed. The design process focused on minimizing network requests by aggregating and page-serving timeline data.

\subsection{comunidade-service Service}
The comunidade-service service acts as the logical coordinator of user social relationships, managing followers and profiles to reduce spatial isolation and support study groups.

\section{Service of Communication, Support and Emergency Response}
This section details the modules focused on immediate student assistance and support communication channels. The chat, call, and emergency services cooperate to ensure that panic alerts or support dialogs are processed safely.

\subsection{chat-service Service}
The chat-service service stores real-time conversation history. The design process aimed to isolate high-write text operations from static tables, preventing write locks on critical user or news tables.

\subsection{emergency-service Service}
The emergency-service service records silent distress alerts containing GPS coordinates, notifying responders immediately. The design prioritized robustness and speed.

\subsection{emergencyCall-service Service}
The emergencyCall-service service handles emergency phone requests. The design provided voice-based redundancy if internet connectivity is unstable.

\subsection{Emergency Services Separation Rationale}
The mental decision to separate chat, emergency, and call services was based on differing latency, availability, and traffic profiles. The emergency service processes lightweight, critical GPS payloads that require absolute transmission reliability. The chat service handles high-frequency conversation text that consumes network bandwidth. Isolating them ensures that conversational traffic spikes do not delay distress alerts. The call service handles telephone integration asynchronously, avoiding database lockups on active voice call requests.

\subsection{ticket-service Service}
The ticket-service service processes user issues as tickets, queueing and routing them to Academic or Frontdesk teams to track resolution statistics.

\subsection{notification-service Service}
The notification-service service distributes real-time notifications to user roles, pushing events directly to the frontend to avoid database polling.

\section{Geolocation, Maps and Routing}
This section explains the spatial integration of NeoLAAC within the city of Covilhã. The services process map coordinates, calculate routes, and render visual emergency geolocalizations.

\subsection{map-service Service}
The map-service service handles campus navigation, computing distances via the Haversine formula and pedestrian road routing using OSRM. It recommends public bus routes when destinations are distant.

\subsection{map-emergency-service Service}
The map-emergency-service service integrates map data with emergency locations on the responder dashboard, mapping the student's visual coordinates in real-time.

\section{Database and Network Infrastructure}
This section describes the database persistence layer and the networking infrastructure needed to run the application securely. It covers database management, backups, secure tunnels, and SSL certificates.

\subsection{DB-central Service}
The DB-central service is the central MariaDB instance, hosting users, tickets, and emergencies to ensure relational integrity.

\subsection{backup-service Service}
The backup-service service runs automated nightly database dumps at 03:00, saving timestamped SQL files to a persistent volume.

The Python code below demonstrates the automatic dump logic of the backup service.

\begin{lstlisting}[language=Python]
# Scheduled backup routine
def run_backup():
    timestamp = datetime.now().strftime("%Y%m%d_%H%M%S")
    main_file = f"{BACKUP_PATH}/main_db_{timestamp}.sql"
    try:
        command = ["mariadb-dump", f"--host={DB_HOST}", f"--user={DB_USER}", f"--password={DB_PASS}", DB_NAME]
        with open(main_file, "w") as f:
            subprocess.run(command, stdout=f, check=True)
    except Exception as e:
        print(f"Error in main backup: {e}")

schedule.every().day.at("03:00").do(run_backup)
while True:
    schedule.run_pending()
    time.sleep(60)
\end{lstlisting}

\subsection{tunnel-service Service}
The tunnel-service service uses the Cloudflare Tunnel agent to expose the local application securely, bypassing dynamic IP and NAT restrictions.

The YAML definition of the tunnel container inside docker-compose is shown below.

\begin{lstlisting}[language=HTML]
# Definition of the Cloudflare Tunnel container
services:
  tunnel:
    image: cloudflare/cloudflared:latest
    container_name: laac-tunnel
    restart: always
    command: tunnel --no-autoupdate run --token ${CLOUDFLARE_TUNNEL_TOKEN}
    depends_on:
      - frontend
\end{lstlisting}

\subsection{certbot Service}
The certbot service automates Let's Encrypt SSL/TLS certificates via ACME HTTP-01 challenges served by Nginx.

\subsection{frontend-service Service}
The frontend-service service serves static UI assets (HTML, CSS, JS) via Nginx and communicates assynchronously with the gateway.

\subsection{dev-service Service}
The dev-service service aggregates local environment variables (.env) and docker compose orchestration to simplify the development setup.

\subsection{logging-service Service}
The logging-service service aggregates audit logs from all microservices, providing a centralized repository for security forensics.

\section{Quality Control and Testing Pipeline}
This section explains the quality control process designed to ensure API route reliability. The testing pipeline operates in an ephemeral environment to validate code without affecting production.

\subsection{unit-tests Service}
The unit-tests service validates gateway endpoints and database integrations using mocks to prevent data modification.

\subsection{pipeline.py Orchestration}
The pipeline.py script runs ephemeral tests, launching temporary container environments, running tests, writing markdown reports, and exiting immediately.

The following excerpt demonstrates the ephemeral execution process of the pipeline script.

\begin{lstlisting}[language=Python]
# Ephemeral test execution routine
def run_pipeline_tests():
    # Run tests in a disposable container environment
    test_result = subprocess.run(
        ["docker-compose", "run", "--rm", "orq", "python", "-m", "pytest", "tests"],
        capture_output=True, text=True
    )
    
    # Write the Markdown report
    report_content = f"# Test Report\n"
    if test_result.returncode == 0:
        report_content += "Status: PASS\n"
    else:
        report_content += f"Status: FAIL\nDetails:\n{test_result.stdout}\n"
        
    with open("test_report.md", "w") as f:
        f.write(report_content)
        
    # Exit and return status code immediately
    sys.exit(test_result.returncode)
\end{lstlisting}

\end{document}
